% Hexagon note, 21 Nov 2026. CC0 1.0.
% Proof: GPT-6 Astra (ChatGPT Pro mode), run HODGE_SIGN, 21 Nov 2026, writing out a sketch of L. Alpoge.
% Text (version 5): Claude Fable 5.2 at effort max, using the text of version 4 by Claude Opus 5.5 at effort max.
% Checks: Claude Opus 5.5 at effort high; Claude Fable 5.2 at effort max (independent re-proof).
\documentclass[11pt,a4paper]{article}
\usepackage[T1]{fontenc}
\usepackage{lmodern}
\usepackage{microtype}
\usepackage[margin=2.7cm]{geometry}
\usepackage{amsmath,amssymb,amsthm}
\PassOptionsToPackage{hyphens}{url}
\usepackage[hidelinks]{hyperref}
\emergencystretch=2em

\theoremstyle{plain}
\newtheorem{theorem}{Theorem}
\newtheorem{lemma}[theorem]{Lemma}
\newtheorem{proposition}[theorem]{Proposition}
\theoremstyle{remark}
\newtheorem{example}[theorem]{Example}
\newtheorem{remark}[theorem]{Remark}

\newcommand{\Fq}{\mathbf F_q}
\newcommand{\Fqb}{\overline{\mathbf F}_q}
\newcommand{\Ql}{\overline{\mathbf Q}_\ell}
\newcommand{\A}{\mathbf A}
\newcommand{\F}{\mathcal F}
\newcommand{\G}{\mathcal G}
\newcommand{\Tr}{\operatorname{Tr}}

\title{A short proof of Deligne's bound for Frobenius eigenvalues,\\ after Alp\"oge}
\author{Entered at deposit}
\date{21 November 2026 \quad CC0 1.0}

\begin{document}
\maketitle

\begin{abstract}
\noindent
We write out in full the proof, sketched by L.~Alp\"oge in a footnote of his note of 20 November 2026, of Deligne's
bound for the eigenvalues of Frobenius. Let $X$ be a separated scheme of finite type over a finite field with $q$
elements, and let $\mathcal F$ be a constructible $\ell$-adic sheaf on $X$ whose Frobenius eigenvalues at each closed
point $x$ have absolute value at most $|\kappa(x)|^{w/2}$. Then the eigenvalues of Frobenius on $H^i_c(X,\mathcal F)$
have absolute value at most $q^{(w+i)/2}$, and they have absolute value exactly $q^{(w+i)/2}$ if $X$ is smooth and
proper and $\mathcal F$ is lisse and pure of weight $w$. The proof is short. Granted the formalism of \'etale
cohomology (the trace formula, the K\"unneth formula, Poincar\'e duality and Artin's vanishing theorem) and a weak
Lefschetz theorem for generic hyperplane sections of affine varieties, it takes less than six pages, and its only
estimate is the Cauchy--Schwarz inequality for a character sum over a finite field. Lefschetz pencils, monodromy and
Rankin's method are not used. The proof was produced by GPT-6 Astra (ChatGPT Pro mode) from that sketch. It was
checked by Claude Opus 5.5 at effort high and, independently, by Claude Fable 5.2 at effort max, which also wrote
this text. No human has checked it.
\end{abstract}

\section{The theorem}

Let $\Fq$ be a finite field, $\ell$ a prime not dividing $q$, and $\iota:\Ql\hookrightarrow\mathbf C$ an embedding of
fields, fixed throughout. The absolute value of an element of $\Ql$ always means the absolute value of its image
under $\iota$. Let $X$ be a separated scheme of finite type over $\Fq$ and $\F$ a constructible $\Ql$-sheaf on $X$.
The groups $H^i_c(X,\F)$ are those of $X\otimes\Fqb$; they are finite-dimensional, and the geometric Frobenius $F$
acts on them. At a closed point $x$ of $X$, with residue field $\kappa(x)$, the geometric Frobenius of $\kappa(x)$
acts on the stalk $\F_{\bar x}$. For a real number $w$, we say that $\F$ \emph{has weights at most $w$} if at every
closed point $x$ all eigenvalues of this local Frobenius have absolute value at most $|\kappa(x)|^{w/2}$, and that
$\F$ is \emph{pure of weight $w$} if they all have absolute value exactly $|\kappa(x)|^{w/2}$.

\begin{theorem}\label{thm:main}
Let $X$ be a separated scheme of finite type over $\Fq$ and $\F$ a constructible $\Ql$-sheaf on $X$.
\begin{enumerate}
\item[(a)] If $\F$ has weights at most $w$, then for every $i$ every eigenvalue of $F$ on $H^i_c(X,\F)$ has absolute
value at most $q^{(w+i)/2}$.
\item[(b)] If $X$ is smooth and proper and $\F$ is lisse and pure of weight $w$, then for every $i$ every eigenvalue
of $F$ on $H^i(X,\F)$ has absolute value $q^{(w+i)/2}$.
\end{enumerate}
\end{theorem}

Part (a) says that the passage from the stalks of $\F$ to its cohomology in degree $i$ raises the weights by at most
$i$. Part (b) says that on a smooth proper scheme nothing is lost: the cohomology in degree $i$ of a pure sheaf is
pure, of weight $w+i$. For the constant sheaf $\Ql$ on a smooth projective variety $X$ the hypothesis of (b) holds
with $w=0$ for every embedding $\iota$, so each eigenvalue of $F$ on $H^i(X,\Ql)$ has absolute value $q^{i/2}$ under
every embedding. This is the Riemann hypothesis for $X$.

Both statements are theorems of Deligne. Part (b) for the constant sheaf is \cite[Th\'eor\`eme 1.6]{WeilI}, apart
from the assertions made there about integrality and independence of $\ell$, which we do not treat. Part (a) is the
statement about absolute values contained in \cite[Th\'eor\`eme 3.3.1]{WeilII} for a scheme over a finite field. That
theorem concerns a morphism $f:X\to S$ and a mixed sheaf, and says that $R^if_!\F$ is mixed of weights at most $w+i$.
For absolute values the relative statement follows from (a), applied to the fibres of $f$, because the stalks of
$R^if_!\F$ are the groups $H^i_c$ of the fibres. Part (a) asks only for a bound at each closed point, for the one
embedding $\iota$, and does not assume that $\F$ is mixed; this is not a gain in generality, since by a theorem of
L.~Lafforgue \cite[Th\'eor\`eme VII.6 and Corollaire VII.8]{Lafforgue} every constructible $\Ql$-sheaf on $X$ is
$\iota$-mixed. Nothing is proved here about the other parts of Deligne's theory: the weight filtration, the local
monodromy theorem and the hard Lefschetz theorem of \cite{WeilII}, the rationality and the functional equation of
zeta functions, or the semisimplicity of Frobenius.

\section{The proof, its source and its length}

\paragraph{The source.}
The note \cite{Alpoge} of L.~Alp\"oge mentions a proof of the theorems of Weil and Deligne that uses only a sign,
based on an identity $qN=M^k+R$, and sketches it in a footnote \cite[p.~2 and footnote 3]{Alpoge}. In the sketch $N$
is a weighted count of the zeros of the form $x_1x_2+\dots+x_{k-1}x_k$, written as a sum over additive characters, and
$M^k$ is the term of the trivial character. An argument with a hyperplane section gives the contribution of middle
cohomology a definite sign, one degree away from that of $M^k$. The remainder $R$ has square-root size, by Parseval's
identity applied $k$ times. Because of the sign, the inverse roots cannot cancel one another, so the bound for $R$ is
a bound for each of them. In higher dimension vectors replace the scalars $x_i$.

The present note writes this argument out in full. Its ideas are those of the sketch. What is supplied here is the
detail: the reduction to the lowest cohomology group of a smooth affine variety, a finite map to affine space that
turns the count into the character sum, the weak Lefschetz theorem that provides the sign, and the passage to all
degrees.

\paragraph{The length.}
The proof is short, and we say in what sense. We take for granted the formalism of $\ell$-adic cohomology as it stood
before \cite{WeilI}: the trace formula, the K\"unneth formula, Poincar\'e duality, the vanishing of the cohomology of
an affine scheme above its dimension, and the exact sequence attached to an open subscheme and its complement. We
also take for granted one theorem of a less elementary kind: for a generic hyperplane section $Y$ of a smooth
affine variety $Z$, and a lisse sheaf on $Z$, restriction from $Z$ to $Y$ is bijective on cohomology below the
dimension of $Y$ and injective in that degree. The proof of it that we use is Deligne's, printed in the appendix of \cite{KatzACT}. It is a few pages long
and rests on his generic base change theorem and on Artin's vanishing theorem; Section~\ref{sec:lefschetz} recalls
it. With these inputs, the proof of Theorem~\ref{thm:main} occupies Sections~\ref{sec:power} to~\ref{sec:proof},
less than six pages, and its only estimate is the Cauchy--Schwarz inequality, applied to a character sum over a
finite field (Lemma~\ref{lem:parseval}).

For comparison, \cite{WeilI} uses the same formalism together with the theory of Lefschetz pencils, a theorem of
Kazhdan and Margulis on their monodromy, and Rankin's method. In \cite{WeilII} the general statement is reduced to
the case of curves, which is proved by a study of monodromy and by the method of Hadamard and de la
Vall\'ee-Poussin. None of these enters below. Laumon \cite{Laumon} gave another proof of the main theorem of
\cite{WeilII}, by the $\ell$-adic Fourier transform, and Kedlaya \cite{Kedlaya} carried that proof over to $p$-adic
cohomology.

For curves, the estimate for the number of points has an elementary proof, by the method of Stepanov
\cite{Bombieri}. Scholl \cite{Scholl} reduced the Riemann hypothesis for smooth projective varieties to the case of
hypersurfaces, and remarked that an elementary proof of the point-count estimate for hypersurfaces, in the manner of
Stepanov, would be of great interest. The proof below is not of that kind. Its estimate is elementary, but it
concerns one explicit hypersurface in a power of the variety, and the passage from that estimate to the eigenvalues
uses cohomology throughout.

\paragraph{Why the proof works.}
By cutting $X$ into pieces and by induction on the dimension, part (a) comes down to one case
(Section~\ref{sec:proof}): $U$ is smooth and affine of pure dimension $d$, the sheaf $\F$ is lisse on $U$, and the
group is $H^d_c(U,\F)$, the lowest one that can be nonzero. In higher degrees the weak Lefschetz theorem expresses
the cohomology of $U$ through that of a hyperplane section, of smaller dimension. In degree $d$ it gives nothing. The
trace formula alone does not help either: it gives the sum of the local traces of $\F$ over the points of $U$ as an
alternating sum over all degrees, in which large eigenvalues from different degrees could cancel.

The way out is to compare two counts. Let $k\ge2$ and let $Y$ be a hypersurface in the $k$-fold product $U^k$. Over a
finite field with $Q$ elements, let $M$ be the sum of the local traces of $\F$ over the points of $U$, and $N$ the
sum of the local traces of $\F\boxtimes\dots\boxtimes\F$ over the points of $Y$. A hypersurface in general position
should carry one $Q$-th of the total, so we put $R=QN-M^k$ and call it the remainder. Suppose that $Y$ has the weak
Lefschetz property: restriction from $U^k$ to $Y$ is bijective in degrees below $kd-1$ and injective in degree
$kd-1$. By duality, every compactly supported cohomology group of $U^k$ above the lowest degree $kd$ is then the
image, bijectively or surjectively, of a group of $Y$. Written with the trace formula, $R$ is an alternating sum over
the groups of $U^k$ and of $Y$, and along these maps the terms cancel in pairs. Two terms survive. One is the lowest
group of $U^k$, which by the K\"unneth formula is the $k$-th tensor power of $H^d_c(U,\F)$. The other is the kernel
$K$ of the map out of the lowest group of $Y$. The first sits in degree $kd$ and is subtracted; the second sits in
degree $kd-1$ and is added; so they enter $R$ with the same sign (Proposition~\ref{prop:remainder}). Up to that one
sign, $R$ is therefore the trace of Frobenius on an actual vector space, the sum of the tensor power and of $K$, and
not a difference of two traces. This matters because of an elementary fact about power sums
(Lemma~\ref{lem:tensor}): if the traces of the powers of an automorphism grow at most like $B^n$, then every one of
its eigenvalues has absolute value at most $B$. For a difference of two traces this is false, since eigenvalues can
cancel.

It remains to find a hypersurface $Y$ with the weak Lefschetz property for which $R$ is small, and the form of its
equation does this. Choose a finite map $\pi$ from $U$ to the affine space $\A^d$ and a nondegenerate bilinear form
$b$ on $\A^d$, and let the equation of $Y$ be
\[
\sum_{i=1}^{k-1}b\bigl(\pi(u_i),\pi(u_{i+1})\bigr)+\sum_{i=1}^kg_i(u_i)+c=0,
\]
where each $g_i$ is a linear combination of coordinates of $U$. Consecutive factors are coupled through $b$, and
there is no other coupling. Detecting the equation by additive characters writes $QN$ as a sum over the characters
$\psi(h\,\cdot)$ of the field. The term with $h=0$ is exactly $M^k$, so $R$ is the sum of the other terms. For
$h\ne0$ the term is a product of $k-1$ matrices with entries $\psi(h\,b(x,y))$, interlaced with diagonal matrices
that carry the local traces. The first kind are $Q^{d/2}$ times unitary matrices, which is Parseval's identity for
the Fourier transform on a finite vector space; the second kind are bounded by a constant times $Q^{w/2}$. Hence
$|R|$ is at most a constant to the power $k$, times $Q^{k(w+d)/2}$, times $Q^{d/2+1}$ (Lemma~\ref{lem:parseval} and
Section~\ref{sec:middle}). This is square-root size, up to the factor $Q^{d/2+1}$, which does not depend on $k$. That
such a $Y$ can be chosen with the weak Lefschetz property is not obvious, because the bilinear part of its equation
is fixed in advance. But $Y$ is a hyperplane section of $U^k$ for the embedding given by the coordinates together
with the fixed bilinear expression, and a generic hyperplane whose coefficient on that expression is $1$ will do
(Lemma~\ref{lem:good}).

Putting the two halves together, every eigenvalue of Frobenius on the $k$-th tensor power of $H^d_c(U,\F)$, in
particular $\alpha^k$ for an eigenvalue $\alpha$ on $H^d_c(U,\F)$, has absolute value at most
$q^{k(w+d)/2+d/2+1}$. Taking $k$-th roots gives $|\alpha|\le q^{(w+d)/2}\cdot q^{(d/2+1)/k}$, and letting $k$ grow
removes the loss.

\section{Power sums of an actual module}\label{sec:power}

Let $V$ be a finite-dimensional vector space over $\Ql$ with an automorphism $\Phi$. Then
$\Tr(\Phi^n\mid V)=\sum_jm_j\lambda_j^n$, where the $\lambda_j$ are the distinct eigenvalues of $\Phi$ and the
$m_j\ge1$ their algebraic multiplicities. The next lemma is the form in which we use the positivity of the $m_j$. It
allows the automorphism, the auxiliary space and the constants to change with $k$, which will be needed.

\begin{lemma}\label{lem:tensor}
Let $V$ be a finite-dimensional $\Ql$-vector space with an automorphism $F$, and let $a$ be a real number. Suppose
that for infinitely many integers $k\ge1$ there are an integer $e_k\ge1$, a finite-dimensional $\Ql$-vector space
$K_k$ with an automorphism $\Phi_k$, a sign $\varepsilon_k=\pm1$, real numbers $C_k>0$ and $c_k\ge0$, and complex
numbers $R_{k,n}$ such that for all $n\ge1$
\[
\varepsilon_kR_{k,n}=\iota\Tr\bigl(F^{e_kn}\mid V^{\otimes k}\bigr)+\iota\Tr\bigl(\Phi_k^n\mid K_k\bigr)
\quad\text{and}\quad|R_{k,n}|\le C_k\,q^{e_kn(ka+c_k)},
\]
and that $c_k/k$ tends to $0$ along these $k$. Then every eigenvalue $\alpha$ of $F$ on $V$ satisfies
$|\iota\alpha|\le q^a$.
\end{lemma}

\begin{proof}
Fix such a $k$. Let $\mu_1,\dots,\mu_s$ be the images under $\iota$ of the distinct eigenvalues of the automorphism
$F^{e_k}\oplus\Phi_k$ of $V^{\otimes k}\oplus K_k$, and $m_1,\dots,m_s\ge1$ their multiplicities, so that
$\varepsilon_kR_{k,n}=\sum_jm_j\mu_j^n$ for all $n\ge1$. The $\mu_j$ are distinct, because $\iota$ is injective, and
nonzero. If $v\ne0$ and $Fv=\alpha v$, then $v\otimes\dots\otimes v$ is an eigenvector of $F^{e_k}$ on
$V^{\otimes k}$ with eigenvalue $\alpha^{ke_k}$, so $\iota(\alpha)^{ke_k}$ is one of the $\mu_j$.

Put $\rho=q^{-e_k(ka+c_k)}$. By the bound on $R_{k,n}$, the power series $\sum_{n\ge1}\varepsilon_kR_{k,n}z^n$
converges in the disc $|z|<\rho$. Near $z=0$ it is equal to the rational function $\sum_jm_j\mu_jz/(1-\mu_jz)$.
Since the $\mu_j$ are distinct and no $m_j$ is zero, this rational function has a pole at each of the points
$1/\mu_j$. It agrees on the disc with a holomorphic function, so none of these points lies in the disc. Hence
$|\mu_j|\le1/\rho$ for every $j$, and in particular $|\iota\alpha|^{ke_k}\le q^{e_k(ka+c_k)}$, that is,
$|\iota\alpha|\le q^{a+c_k/k}$. Now let $k$ tend to infinity.
\end{proof}

\section{A character sum along a path}\label{sec:parseval}

The estimate of this note is the following. The sum runs over $k$ points of a finite vector space; each point is
weighted by an arbitrary bounded function, and the phase couples only consecutive points.

\begin{lemma}\label{lem:parseval}
Let $E$ be a finite field with $Q$ elements, $\psi$ a nontrivial additive character of $E$, $d\ge1$, $A=E^d$ and
$k\ge2$. Let $b_1,\dots,b_{k-1}$ be nondegenerate bilinear forms on $A$ with values in $E$, and let
$t_1,\dots,t_k:A\to\mathbf C$ be functions with $|t_i(x)|\le T$ for all $i$ and $x$. Then for every $h\in E^\times$
and $c\in E$
\[
\Bigl|\sum_{x\in A^k}t_1(x_1)\cdots t_k(x_k)\,\psi\Bigl(h\Bigl(\sum_{i=1}^{k-1}b_i(x_i,x_{i+1})+c\Bigr)\Bigr)\Bigr|
\le T^kQ^{d(k+1)/2}.
\]
\end{lemma}

\begin{proof}
For $i<k$ let $K_i$ be the matrix with rows and columns indexed by $A$ and entries $K_i(x,y)=\psi(h\,b_i(x,y))$. The
entry of $K_iK_i^*$ at $(x,x')$ is $\sum_{y\in A}\psi\bigl(h\,b_i(x-x',y)\bigr)$. If $x\ne x'$, the map
$y\mapsto h\,b_i(x-x',y)$ is a nonzero linear form on $A$, because $b_i$ is nondegenerate and $h\ne0$, and the sum
of a nontrivial character over $A$ is zero. If $x=x'$ the sum is $Q^d$. So $K_iK_i^*=Q^d$, and $K_i$ is $Q^{d/2}$
times a unitary matrix. Let $D_i$ be the diagonal matrix with entries $t_i(x)$; its norm is at most $T$. Regard
$t_1$ and $t_k$ as vectors; their norms are at most $TQ^{d/2}$. The sum in the lemma equals
$\psi(hc)\,t_1^{\mathsf T}K_1D_2K_2\cdots D_{k-1}K_{k-1}t_k$, and by the Cauchy--Schwarz inequality its absolute value
is at most
\[
\|t_1\|\,\|t_k\|\,\prod_{i=1}^{k-1}\|K_i\|\,\prod_{i=2}^{k-1}\|D_i\|\le(TQ^{d/2})^2\,Q^{d(k-1)/2}\,T^{k-2}
=T^kQ^{d(k+1)/2}.\qedhere
\]
\end{proof}

The trivial bound for the sum is $T^kQ^{dk}$. The lemma gains a factor $Q^{d/2}$ for each of the $k-1$ couplings.

\section{The remainder of a point count}\label{sec:remainder}

We recall what duality gives. Let $Z$ be a smooth scheme of pure dimension $m$ over a finite field $\mathbf F_r$ and
$\G$ a lisse $\Ql$-sheaf on $Z$. Poincar\'e duality \cite[XVIII, \S3]{SGA4} is a perfect pairing, compatible with
Frobenius, between $H^i_c(Z,\G)$ and $H^{2m-i}(Z,\G^\vee)(m)$. Here $(m)$ is the Tate twist; on a space twisted by
$(-1)$ the eigenvalues of Frobenius are multiplied by $r$. If $Z$ is affine, then $H^i(Z,\G^\vee)=0$ for $i>m$ by
Artin's theorem \cite[XIV, 3.1]{SGA4}, and therefore $H^i_c(Z,\G)=0$ for $i<m$. If $Y\subset Z$ is a smooth closed
subscheme of pure dimension $m-1$, we use duality on $Z$ and on $Y$: the transpose of the restriction map
$H^a(Z,\G^\vee)\to H^a(Y,\G^\vee)$ is a map
\[
H^j_c(Y,\G)(-1)\longrightarrow H^{j+2}_c(Z,\G),\qquad j=2m-2-a,
\]
compatible with Frobenius. We call it the \emph{Gysin map}. It is bijective when the restriction map is bijective,
and surjective when the restriction map is injective.

\begin{proposition}\label{prop:remainder}
Let $U$ be a smooth affine scheme of pure dimension $d\ge1$ over a finite field $\mathbf F_r$, and $\F$ a lisse
$\Ql$-sheaf on $U$. Let $k\ge1$, $X=U^k$, $D=kd$ and $\F_k=\F\boxtimes\dots\boxtimes\F$ on $X$. Let $Y\subset X$ be a
smooth closed subscheme of pure dimension $D-1$, possibly empty, defined over $\mathbf F_r$, such that over
$\overline{\mathbf F}_r$ the restriction map
\[
H^a(X,\F_k^\vee)\longrightarrow H^a(Y,\F_k^\vee)
\]
is an isomorphism for $a<D-1$ and injective for $a=D-1$. For $n\ge1$ let $M_n$ be the sum of the traces
$\Tr(F_{u,n}\mid\F_{\bar u})$ over $u\in U(\mathbf F_{r^n})$, let $N_n$ be the corresponding sum for $\F_k$ over
$Y(\mathbf F_{r^n})$, and put $R_n=r^nN_n-M_n^k$. Let $F$ be the geometric Frobenius of $\mathbf F_r$, and let $K$ be
the kernel of the Gysin map $H^{D-1}_c(Y,\F_k)(-1)\to H^{D+1}_c(X,\F_k)$. Then for all $n\ge1$
\[
(-1)^{D-1}R_n=\Tr\bigl(F^n\mid H^d_c(U,\F)^{\otimes k}\bigr)+\Tr\bigl(F^n\mid K\bigr).
\]
\end{proposition}

\begin{proof}
The schemes $U$, $X$ and $Y$ are smooth and affine, of pure dimensions $d$, $D$ and $D-1$. By what was recalled, the
groups $H^j_c(U,\F)$ vanish for $j<d$, the groups $H^j_c(X,\F_k)$ for $j<D$, and the groups $H^j_c(Y,\F_k)$ for
$j<D-1$. The hypothesis on restriction says that the Gysin map $H^j_c(Y,\F_k)(-1)\to H^{j+2}_c(X,\F_k)$ is an
isomorphism for $j>D-1$ and surjective for $j=D-1$.

The stalk of $\F_k$ at a point $(u_1,\dots,u_k)$ is the tensor product of the stalks of $\F$ at the $u_i$, so the
sum of the traces of $\F_k$ over $X(\mathbf F_{r^n})=U(\mathbf F_{r^n})^k$ is $M_n^k$. By the trace formula
\cite[Rapport, Th\'eor\`eme 3.2]{SGA4h}, applied on $X$ and on $Y$,
\begin{align*}
M_n^k&=\sum_{b\ge0}(-1)^{D+b}\,\Tr\bigl(F^n\mid H^{D+b}_c(X,\F_k)\bigr),\\
r^nN_n&=\sum_{a\ge0}(-1)^{D-1+a}\,\Tr\bigl(F^n\mid H^{D-1+a}_c(Y,\F_k)(-1)\bigr).
\end{align*}
Therefore, in $(-1)^{D-1}R_n$, the trace on $H^{D-1+a}_c(Y,\F_k)(-1)$ enters with the sign $(-1)^a$ and the
trace on $H^{D+b}_c(X,\F_k)$ with the sign $(-1)^b$. For $a\ge1$ the Gysin map identifies $H^{D-1+a}_c(Y,\F_k)(-1)$ with
$H^{D+1+a}_c(X,\F_k)$, which enters with the opposite sign, and the two traces cancel. Three terms remain: the trace
on $H^{D-1}_c(Y,\F_k)(-1)$, minus the trace on $H^{D+1}_c(X,\F_k)$, plus the trace on $H^D_c(X,\F_k)$. The second
space is the quotient of the first by $K$, so the first two terms together are the trace on $K$. Finally, by the
K\"unneth formula \cite[XVII, 5.4.3]{SGA4} the group $H^D_c(X,\F_k)$ is the sum of the tensor products
$H^{j_1}_c(U,\F)\otimes\dots\otimes H^{j_k}_c(U,\F)$ with $j_1+\dots+j_k=D$. Since the factors vanish for $j_i<d$,
only $j_1=\dots=j_k=d$ contributes, and $H^D_c(X,\F_k)=H^d_c(U,\F)^{\otimes k}$.
\end{proof}

Nothing is assumed about the higher cohomology groups of $U$ or of $Y$: they cancel through actual isomorphisms, and
the surjection at the border leaves an actual kernel.

\begin{example}
Let $U=\A^1$, $\F=\Ql$ and $k=2$, and let $Y\subset\A^2$ be the conic $xy+ax+by+c=0$ with $ab\ne c$. It is
isomorphic to the multiplicative group, so $N_n=r^n-1$, while $M_n=r^n$; hence $R_n=-r^n$ and
$(-1)^{D-1}R_n=r^n$. On the other side $H^1_c(\A^1,\Ql)=0$, and $K=H^1_c(Y,\Ql)(-1)$ has dimension one, with
Frobenius acting by $r$. The group $H^2_c(Y,\Ql)(-1)$ is identified with $H^4_c(\A^2,\Ql)$ and has cancelled.
\end{example}

\begin{remark}
The space $K$ is the analogue, for compact supports on an affine variety, of the part of the cohomology of a
hyperplane section that the Gysin map kills. In Grothendieck's notes on the standard conjectures in the Montpellier
archive this part is his \emph{partie \'evanescente} $E^{n-1}$ (cote 16, scan 40, in the transcription
\cite{Archive}). The comparison plays no role in the proof.
\end{remark}

\section{Hypersurfaces with the weak Lefschetz property}\label{sec:lefschetz}

\begin{theorem}[Deligne; {\cite[Appendix, A.5, with A.3 and A.4]{KatzACT}}]\label{thm:deligne}
Let $Z$ be a separated scheme of finite type over an algebraically closed field of characteristic different from
$\ell$, let $\varphi:Z\to\mathbf P^N$ be a quasi-finite morphism, and let $\mathcal K$ be a perverse sheaf on $Z$.
Then there is a dense open set of hyperplanes $H\subset\mathbf P^N$ for which the restriction map
\[
H^a(Z,\mathcal K)\longrightarrow H^a(\varphi^{-1}H,\mathcal K)
\]
is an isomorphism for $a<-1$ and injective for $a=-1$.
\end{theorem}

Katz attributes this proof to Deligne \cite[p.~152]{KatzACT}. It runs as follows. Let
$\alpha:Z\smallsetminus\varphi^{-1}H\to Z$ be the inclusion. The assertion is equivalent to the vanishing of
$H^a(Z,\alpha_!\alpha^*\mathcal K)$ for $a<0$. For generic $H$, the generic base change theorem
\cite[Th\'eor\`emes de finitude, 1.9]{SGA4h}, applied over the space of hyperplanes, identifies these groups with
the compactly supported cohomology of the affine scheme $\mathbf P^N\smallsetminus H$ with coefficients in a direct
image of $\alpha^*\mathcal K$; duality and Artin's theorem then give the vanishing. No weights occur. If $Z$ is
smooth of pure dimension $D$ and $\G$ is a lisse sheaf on $Z$, then $\G[D]$ is perverse, and the conclusion reads:
the restriction map $H^a(Z,\G)\to H^a(\varphi^{-1}H,\G)$ is an isomorphism for $a<D-1$ and injective for $a=D-1$.

We need this for hypersurfaces whose equation has a fixed part. Corollary 3.4.1 of \cite{KatzACT} states exactly that
form, but it is deduced there from the main theorem of that article, whose proof is by the Fourier transform. The
next lemma deduces it from Theorem~\ref{thm:deligne} instead, so that only the inputs named above are used; it also
gives the smoothness of the hypersurface.

\begin{lemma}\label{lem:good}
Let $V$ be a smooth affine scheme of pure dimension $D\ge1$ over $\Fqb$, let $y_1,\dots,y_M$ be functions on $V$
such that $\Gamma=(y_1,\dots,y_M):V\to\A^M$ is a closed immersion, and let $\G$ be a lisse $\Ql$-sheaf on $V$.
Consider the hypersurfaces
\[
Y_{\lambda,\mu}=V\cap\Bigl(\,\sum_{m=1}^M\lambda_my_m+\mu=0\Bigr),\qquad\lambda\ne0.
\]
For $(\lambda,\mu)$ in a dense open subset of the space of all parameters, and also for $(\lambda,\mu)$ in a dense
open subset of the space of parameters with $\lambda_M=1$, the scheme $Y_{\lambda,\mu}$ is smooth, of pure dimension
$D-1$ or empty on each connected component of $V$, and the restriction map
$H^a(V,\G)\to H^a(Y_{\lambda,\mu},\G)$ is an isomorphism for $a<D-1$ and injective for $a=D-1$.
\end{lemma}

\begin{proof}
Regard $\A^M$ as the complement of the hyperplane at infinity in $\mathbf P^M$. The hyperplanes of $\mathbf P^M$
other than the one at infinity are the zero sets of the equations $\sum_m\lambda_my_m+\mu=0$ with $\lambda\ne0$, and
the inverse image of such a hyperplane under $\Gamma$ is $Y_{\lambda,\mu}$.

The restriction property holds for a dense open set of hyperplanes by Theorem~\ref{thm:deligne}, applied to
$\varphi=\Gamma$, which is quasi-finite, and to the perverse sheaf $\G[D]$.

Smoothness holds for a dense open set of hyperplanes by a count of dimensions. Consider the pairs $(v,H)$ of a point
$v$ of $V$ and a hyperplane $H$ that contains the tangent space of $\Gamma(V)$ at $\Gamma(v)$. Since $\Gamma$ is an
immersion, this tangent space has dimension $D$, and the hyperplanes containing it form a projective space of
dimension $M-D-1$. So the set of such pairs has dimension at most $D+(M-D-1)=M-1$, and its image in the space of
hyperplanes, which has dimension $M$, is not dense. If $H$ lies outside the closure of this image, the equation of
$H$ restricts to a nonzero linear form on the tangent space of $V$ at each point of $\Gamma^{-1}H$, so
$\Gamma^{-1}H$ is smooth, and of pure dimension $D-1$ or empty on each component of $V$. This argument is valid in
every characteristic, because $\Gamma$ is an immersion; compare \cite[I.6.10]{Jouanolou}.

The hyperplanes with both properties form a dense open subset of the space of all hyperplanes. It meets the open
subset $\lambda_M\ne0$ in a dense open subset, and on that subset we may scale the equation so that $\lambda_M=1$.
\end{proof}

\section{The lowest degree of a smooth affine variety}\label{sec:middle}

\begin{theorem}\label{thm:middle}
Let $U$ be a smooth affine scheme of pure dimension $d\ge1$ over $\Fq$ and $\F$ a lisse $\Ql$-sheaf on $U$. Suppose
that there are a constant $C>0$ and a real number $w$ with
\[
|\iota\Tr(F_{u,n}\mid\F_{\bar u})|\le C\,q^{nw/2}\qquad\text{for all }n\ge1\text{ and }u\in U(\mathbf F_{q^n}).
\]
Then every eigenvalue of $F$ on $H^d_c(U,\F)$ has absolute value at most $q^{(w+d)/2}$.
\end{theorem}

\begin{proof}
\emph{Choices.} By Noether normalization \cite[Tag 00OY]{Stacks} there is a finite morphism $\pi:U\to\A^d$ over
$\Fq$. Let $B$ be the number of elements in a set of generators of the coordinate ring $\mathcal O(U)$ as a module
over that of $\A^d$. Over any field, each fibre of $\pi$ is the spectrum of an algebra of dimension at most $B$, so
it has at most $B$ rational points. Choose generators $z_1,\dots,z_s$ of $\mathcal O(U)$ as an $\Fq$-algebra, so that
$z=(z_1,\dots,z_s):U\to\A^s$ is a closed immersion, and a nondegenerate bilinear form $b$ on $\A^d$ defined over
$\Fq$. Fix $k\ge2$ and put $D=kd$.

\emph{The hypersurface.} For parameters $a=(a_{ij})$, with $1\le i\le k$ and $1\le j\le s$, and $c$, let $P_{a,c}$
be the function on $U^k$
\[
P_{a,c}(u_1,\dots,u_k)=\sum_{i=1}^{k-1}b\bigl(\pi(u_i),\pi(u_{i+1})\bigr)+\sum_{i=1}^k\sum_{j=1}^sa_{ij}\,z_j(u_i)+c,
\]
and let $Y_{a,c}\subset U^k$ be its zero set. Over $\Fqb$ the map
\[
\Gamma:\ u\longmapsto\Bigl(z(u_1),\dots,z(u_k),\ \sum_{i=1}^{k-1}b\bigl(\pi(u_i),\pi(u_{i+1})\bigr)\Bigr)
\]
is a closed immersion of $U^k$ into $\A^{sk+1}$, because its first $sk$ coordinates already are one. The $Y_{a,c}$
are exactly the hypersurfaces of Lemma~\ref{lem:good} for $\Gamma$ whose coefficient on the last coordinate is $1$.
By that lemma, applied to the sheaf $\F_k^\vee$, there is a dense open set of parameters $(a,c)$ for which $Y_{a,c}$
is smooth of pure dimension $D-1$, or empty, and has the restriction property required in
Proposition~\ref{prop:remainder}. A dense open subset of an affine space over $\Fqb$ has a point with coordinates in
a finite field. Fix such a point $(a,c)$, with coordinates in $\mathbf F_r$, $r=q^e$, and write $P=P_{a,c}$ and
$Y=Y_{a,c}$. The exponent $e$ and the point $(a,c)$ depend on $k$; this is allowed in Lemma~\ref{lem:tensor}.

\emph{The remainder as a character sum.} Let $n\ge1$, $Q=r^n$ and $E=\mathbf F_Q$, and let $M_n$, $N_n$ and
$R_n=QN_n-M_n^k$ be formed over $\mathbf F_r$ as in Proposition~\ref{prop:remainder}. Let $\psi$ be a nontrivial
additive character of $E$, and for $u\in U(E)$ put $\tau(u)=\iota\Tr(F_{u,en}\mid\F_{\bar u})$; by hypothesis
$|\tau(u)|\le CQ^{w/2}$. For a point $u=(u_1,\dots,u_k)$ of $U^k$ with values in $E$, the sum
$\sum_{h\in E}\psi(hP(u))$ is $Q$ if $P(u)=0$ and $0$ otherwise. Therefore
\[
Q\cdot\iota N_n=\sum_{h\in E}\ \sum_{u\in U(E)^k}\tau(u_1)\cdots\tau(u_k)\,\psi\bigl(hP(u)\bigr).
\]
In the inner sum we collect the $u_i$ according to their images $x_i=\pi(u_i)$ in $E^d$. With
\[
t_{i,h}(x)=\sum_{u\in U(E),\ \pi(u)=x}\tau(u)\,\psi\Bigl(h\sum_{j=1}^sa_{ij}z_j(u)\Bigr),
\]
the inner sum becomes
\[
\psi(hc)\sum_{x\in(E^d)^k}t_{1,h}(x_1)\cdots t_{k,h}(x_k)\,\psi\Bigl(h\sum_{i=1}^{k-1}b(x_i,x_{i+1})\Bigr).
\]
For $h=0$ all phases are $1$, and the inner sum is $\bigl(\sum_{u\in U(E)}\tau(u)\bigr)^k=(\iota M_n)^k$. For
$h\ne0$, each $t_{i,h}(x)$ is a sum of at most $B$ terms of absolute value at most $CQ^{w/2}$, so
Lemma~\ref{lem:parseval} applies with $T=BCQ^{w/2}$ and all $b_i$ equal to $b$. There are $Q-1$ such $h$, and we
obtain
\[
|\iota R_n|=\bigl|Q\cdot\iota N_n-(\iota M_n)^k\bigr|\le(Q-1)\,(BC)^kQ^{kw/2}\,Q^{d(k+1)/2}
\le(BC)^k\,q^{en\{k(w+d)/2+d/2+1\}}.
\]

\emph{Conclusion.} The geometric Frobenius of $\mathbf F_r$ is $F^e$. By Proposition~\ref{prop:remainder}, applied
over $\mathbf F_r$,
\[
(-1)^{D-1}\,\iota R_n=\iota\Tr\bigl(F^{en}\mid H^d_c(U,\F)^{\otimes k}\bigr)+\iota\Tr\bigl((F^e)^n\mid K\bigr)
\]
for all $n\ge1$. These are the hypotheses of Lemma~\ref{lem:tensor} for $V=H^d_c(U,\F)$ and $a=(w+d)/2$, with
$e_k=e$, $K_k=K$, $\varepsilon_k=(-1)^{D-1}$, $C_k=(BC)^k$ and $c_k=d/2+1$, for every $k\ge2$. The lemma gives the
theorem.
\end{proof}

\section{Proof of Theorem~\ref{thm:main}}\label{sec:proof}

A bound for the eigenvalues may be proved after replacing $\Fq$ by a finite extension $\mathbf F_{q^e}$: the
eigenvalues of $F^e$ are the $e$-th powers of those of $F$, so a bound $q^{e(w+i)/2}$ for the former gives the bound
$q^{(w+i)/2}$ for the latter.

\medskip\noindent\emph{Part (a), for $X$ smooth and affine of pure dimension $d$ and $\F$ lisse.} We prove the
statement by induction on $d$, for all finite fields at once. Let $d=0$. Then $H^0_c(X,\F)$ is the sum of the stalks
of $\F$ at the geometric points of $X$. If $x$ is a closed point of degree $m$, the Frobenius $F$ permutes the $m$
geometric points above $x$ cyclically, and $F^m$ acts on each of their stalks as the local Frobenius at $x$. So for
an eigenvalue $\alpha$ of $F$ on the part of $H^0_c$ above $x$, $\alpha^m$ is an eigenvalue of the local Frobenius,
and $|\alpha|^m\le q^{mw/2}$.

Let $d\ge1$. For $i<d$ the group $H^i_c(X,\F)$ vanishes (Section~\ref{sec:remainder}). For $i=d$ we apply
Theorem~\ref{thm:middle}, with $C$ the rank of $\F$: if $u$ is a point with values in $\mathbf F_{q^n}$ over a closed
point $x$ of degree $m$, then $F_{u,n}$ is the $(n/m)$-th power of the local Frobenius at $x$, so its eigenvalues
have absolute value at most $q^{nw/2}$. For $i>d$ we cut by a hyperplane. Choose a closed immersion of $X$ into an
affine space. By Lemma~\ref{lem:good}, applied to this immersion, to all hyperplanes and to the sheaf $\F^\vee$,
there is, over some finite extension $\mathbf F_{q^{e}}$, a hyperplane section $Y$ of $X$ that is smooth and affine
of pure dimension $d-1$, or empty, and for which restriction from $X$ to $Y$ on the cohomology of $\F^\vee$ is
bijective below degree $d-1$ and injective in degree $d-1$. Dually (Section~\ref{sec:remainder}) the Gysin map
\[
H^{i-2}_c(Y,\F)(-1)\longrightarrow H^i_c(X,\F)
\]
is surjective for $i=d+1$ and bijective for $i>d+1$. By the induction hypothesis for $Y$ over $\mathbf F_{q^e}$, the
eigenvalues of $F^e$ on $H^{i-2}_c(Y,\F)$ have absolute value at most $q^{e(w+i-2)/2}$. The twist multiplies them by
$q^e$. So the eigenvalues of $F^e$ on $H^i_c(X,\F)$ have absolute value at most $q^{e(w+i)/2}$, which is the bound.

\medskip\noindent\emph{Part (a) in general.} The cohomology of $X$ is that of the reduced scheme, so we may assume
$X$ reduced. We use induction on the dimension of $X$; in dimension $0$ the scheme $X$ is smooth and affine and
$\F$ is lisse, which is the case already treated. Since $\Fq$ is perfect, $X$ has a dense open subscheme $U$ that is
smooth and affine, of pure dimension on each of its connected components, and on which $\F$ is lisse. Its complement
$Z$ has smaller dimension, and the restriction of $\F$ to $Z$ again has weights at most $w$. The sequence
\[
H^i_c(U,\F)\longrightarrow H^i_c(X,\F)\longrightarrow H^i_c(Z,\F)
\]
is exact and compatible with Frobenius; it is part of the long exact sequence of
$0\to j_!j^*\F\to\F\to i_*i^*\F\to0$, where $j$ and $i$ are the inclusions of $U$ and $Z$. So every eigenvalue of
$F$ on $H^i_c(X,\F)$ is an eigenvalue on $H^i_c(U,\F)$ or on $H^i_c(Z,\F)$. The first are bounded by the case
already treated, applied to the components of $U$, and the second by the induction hypothesis.

\medskip\noindent\emph{Part (b).} Treating the connected components of $X$ separately, we may assume that $X$ has
pure dimension $d$. Since $X$ is proper, $H^i(X,\F)=H^i_c(X,\F)$, and part (a) gives the upper bound. The dual sheaf
$\F^\vee$ is lisse and pure of weight $-w$, because the eigenvalues of a local Frobenius on its stalk are the
inverses of those on the stalk of $\F$; here we use that $\F$ is pure at every point, and not only bounded.
Poincar\'e duality is a perfect pairing of $H^i(X,\F)$ with $H^{2d-i}(X,\F^\vee)$ into $\Ql(-d)$, compatible with
Frobenius. So if $\alpha$ is an eigenvalue of $F$ on $H^i(X,\F)$, then $q^d/\alpha$ is an eigenvalue of $F$ on
$H^{2d-i}(X,\F^\vee)$. Part (a), applied to $\F^\vee$, gives $|q^d/\alpha|\le q^{(-w+2d-i)/2}$, that is,
$|\alpha|\ge q^{(w+i)/2}$. \qed

\section*{Provenance and checks}

The proof is Part I, Section 1 of the output of a research run of GPT-6 Astra (ChatGPT Pro mode) on 21 November 2026,
whose task was to work out footnote 3 of \cite{Alpoge}. The same day Claude Opus 5.5 at effort high checked every
step of it. Claude Fable 5.2 at effort max, without seeing that check, re-proved every step and found no gap; it read
the pages of \cite{KatzACT} used here (pp.~152--158, 179--181 and 219--222) and re-proved
Theorem~\ref{thm:deligne} independently. One correction came from that second check and is adopted here: the weak
Lefschetz input is taken from Deligne's proof in Katz's appendix, through Lemma~\ref{lem:good}, and not from the
printed proof of Katz's Corollary 3.4.1.

The other references were not reopened for this text. The foundational ones are cited with the locators of the run.
What the first two sections say about \cite{Bombieri}, \cite{WeilI}, \cite{WeilII}, \cite{Kedlaya},
\cite{Lafforgue} and \cite{Laumon} is from memory, and what they say about \cite{Scholl} rests on a reading of its introduction by
Claude Opus 5.5 at effort max. No human has checked the proof or this text.

The prompt, the output of the run and the checks are public at

{\raggedright\smallskip\noindent\url{https://github.com/KokunoYumeto/grothendieck-montpellier-archive/tree/main/research/grh_20261120/deligne_sign_only}.\par}

\smallskip
This text was written by Claude (Claude Fable 5.2 at effort max), using an earlier text by Claude Opus 5.5 at effort
max. It is dedicated to the public domain (CC0 1.0).

\begin{thebibliography}{99}
\bibitem{Alpoge} L.~Alp\"oge, note on the Riemann hypothesis, 20 November 2026, \url{https://alpo.ge/remark.pdf}.
\bibitem{Bombieri} E.~Bombieri, Counting points on curves over finite fields (d'apr\`es S.~A.~Stepanov), S\'eminaire
Bourbaki 1972/73, Expos\'e 430, Lecture Notes in Math. 383, Springer, 1974, 234--241.
\bibitem{WeilI} P.~Deligne, La conjecture de Weil.~I, Publ. Math. IH\'ES 43 (1974), 273--307.
\bibitem{WeilII} P.~Deligne, La conjecture de Weil.~II, Publ. Math. IH\'ES 52 (1980), 137--252.
\bibitem{Jouanolou} J.-P.~Jouanolou, Th\'eor\`emes de Bertini et applications, Birkh\"auser, 1983.
\bibitem{KatzACT} N.~M.~Katz, Affine cohomological transforms, perversity, and monodromy, J. Amer. Math. Soc. 6
(1993), 149--222.
\bibitem{Kedlaya} K.~S.~Kedlaya, Fourier transforms and $p$-adic ``Weil II'', Compos. Math. 142 (2006), 1426--1450.
\bibitem{Lafforgue} L.~Lafforgue, Chtoucas de Drinfeld et correspondance de Langlands, Invent. Math. 147 (2002),
1--241.
\bibitem{Laumon} G.~Laumon, Transformation de Fourier, constantes d'\'equations fonctionnelles et conjecture de
Weil, Publ. Math. IH\'ES 65 (1987), 131--210.
\bibitem{Scholl} A.~J.~Scholl, Hypersurfaces and the Weil conjectures,
\href{https://arxiv.org/abs/0912.2073}{arXiv:0912.2073}.
\bibitem{SGA4} M.~Artin, A.~Grothendieck, J.-L.~Verdier, Th\'eorie des topos et cohomologie \'etale des sch\'emas
(SGA~4), Lecture Notes in Math. 269, 270, 305, Springer, 1972--1973.
\bibitem{SGA4h} P.~Deligne, Cohomologie \'etale (SGA~$4\frac12$), Lecture Notes in Math. 569, Springer, 1977.
\bibitem{Stacks} The Stacks Project, \url{https://stacks.math.columbia.edu/tag/00OY}.
\bibitem{Archive} Transcription of A.~Grothendieck's papers in the Montpellier archive,
\url{https://github.com/KokunoYumeto/grothendieck-montpellier-archive}.
\end{thebibliography}

\end{document}
